% Opening scope/coverage note for the complete supplementary proof volume.
% Prepared 7 October 2026. Parent supplies longtable and booktabs.
\section{Coverage, dependencies, and scope of the supplementary results}
\label{sec:supplement-coverage}

This supplement supplies the longer proofs and auxiliary results underlying
the main article. It also records how the earlier special cases are covered
by the final statements. All graphs are finite and nonempty, and all second
neighborhoods are defined by directed distance \emph{exactly} two. In
particular, first neighbors are excluded from second neighborhoods. Write
\(g_D(v)=d_D^+(v)-d_D^{++}(v)\); thus a Seymour vertex has
\(g_D(v)\le0\), and an all-deficient graph has positive gap everywhere.

\paragraph{The hypotheses are part of each result.}
The residual arguments use a strongly connected all-deficient graph that
is minimal under deletion of \emph{every nonempty set of arcs}, on its fixed
vertex set. We call this the MIN setting. Single-arc criticality alone is
not the stated hypothesis. Sink-component arguments extend the numerical
missing-pair bound to arc-set-minimal graphs without strong connectivity.
The incoming-pattern certificate theorems and the basic closure theorems
apply to arbitrary oriented graphs and do not require MIN. The class-cover
and cloning normal forms instead use their explicitly specified
minimum-order or lexicographic extremal choices. An intermediate closure
graph, or a graph obtained by vertex deletion, is never assumed to inherit
all-deficiency or minimality without a separate argument.

\paragraph{Incoming-pattern certificates.}
The main article proves the transfer from representative cores, its
tournament case, the at-most-five maximal incoming-class theorem, the
iterated core reduction, and the equivalence between universal pattern
feasibility and the second-neighborhood conjecture. The accompanying
integer tables and independent verifiers certify every core of order at
most five. This bound concerns core order, not the total graph order.

\paragraph{Cloning and extremal forms.}
The main article gives simultaneous original-snapshot row cloning, the
minimum-order whole-class gap-one predecessor cover, and the specified
extremal true-twin normal form. Gap-one vertices here must not be confused
with residual-one vertices. Classical gap-one predecessor and cycle facts
are distinguished from their class-level refinements.

\paragraph{Residual foundation.}
The supplement supplies the external-boundary deletion lemma, target
packet, \(\eta\ge0\), \(\sum\eta=2M-n\), zero-residual protection,
unit residual budgets, and the shared exceptional-partner remainder.
These are the local inputs to the MIN statements in the article.

\paragraph{Completion and residual margin.}
Canonical convenient orientations, witness transfer, safe retention,
two-positive-support exclusion and the fixed heavy-cover criterion lead
to the one-exception theorem. The remaining mixed cases
\((m,2,1)\), \((2,2,2)\), and \((2,2,1,1)\) are proved here in full.
Together they yield the article's parity-sensitive bound on missing
pairs. All comparisons of completed gaps are made in the specified
completion.

\paragraph{Closure and equality.}
The main article develops iterated forced and single weak protection
steps, exact original feed neighborhoods, strict far-head surplus, and
the original partner-cycle and seed-path conclusions. The supplement
retains the one-round antecedents, packet refinements, low-residual
equality shapes, near-whole feed boundary bound, and the separately
scoped additional safe operator.

\paragraph{Additional transfer tools.}
Complete supplementary arguments cover all-keeper and overlap budgets,
multiple retained cross-witness arcs, flexible exception-clique feed
reduction, arbitrary outgoing-star deletion identities, and the
gap-two witness multiplicity consequences. These have their own
local hypotheses and do not assert universal completion existence.

\paragraph{Dense odd boundary.}
At \(n=2\delta+3\), the supplement includes the target-capacity core,
minimum-degree predecessor bounds, even-order and odd-deletion
identities, certified basins, heavy-set capacity, and the degree- and
residual-dependent four-heavy restrictions. These are independent
auxiliary consequences rather than a solution of the dense boundary.

\paragraph{Obstructions and verification.}
Cyclic amplification, exact weighted transport, its dual formulation,
and explicit controls document limits of proposed shortcuts. The
artifact index distinguishes exact integer certificate checking,
finite local tests, solver-dependent infeasibility logs, and heuristic
search outcomes. None of the control graphs is a counterexample to the
second-neighborhood conjecture.

\paragraph{The short dependency chain for the numerical bound.}
In MIN, at least three vertices have positive residual, and the
one-exception theorem requires at least two \emph{actual} bad-direction
failure sources with residual at least two. Consequently the only
positive residual partitions with total at most six are
\[
 (2,2,1),\qquad(3,2,1),\qquad(2,2,2),\qquad(2,2,1,1).
\]
The general \((m,2,1)\) theorem eliminates the first two, and the other
two exclusions eliminate the remainder. Hence
\[
 R=2M-n\ge7,\qquad
 M\ge\left\lceil\frac n2\right\rceil+3+
       \mathbf1_{\{n\ \mathrm{even}\}}.
\]
The even-order residual is at least eight by parity. The sink-component
extension adds the stated penalty for vertices outside a sink component;
equality in the uniform bound forces strong connectivity. These bounds
are assertions about arc-set-minimal counterexamples. Since arc deletion
increases the number of missing pairs, they do not directly prove the
conjecture for arbitrary denser graphs.

\paragraph{Earlier exclusions and their final disposition.}
The unrestricted one-exception theorem excludes all allocations
\((m,1^k)\), for every \(m\ge2\) and finite \(k\), and also covers the
pure-unit situation. It therefore subsumes the earlier finite unit-support
thresholds, the standalone all-unit exclusion, and the cases
\((m,1,1)\), \((m,1,1,1)\), and \((2,1,1,1,1)\).
The complete two-support theorem supersedes its smaller residual and
restricted-order cases. The final numerical bound supersedes the
successive exclusions of residuals zero through five and the older
\(n=19,\delta=8\) conclusions \(M\ge6\) and \(M\ge7\); in that
MIN case the present bound is \(M\ge13\). Useful local lemmas from
these arguments remain available, including full unit packets,
unit--unit witness tails, simultaneous two-vertex priorities, and the
matching-module deletion identities. They are not counted again as
independent improvements of the final bound.

\paragraph{Zero gap and strong-feed equality are distinct branches.}
For the terminal FORCED/WEAK construction, a positive-original-gap feed
always satisfies \(t_D(f)\ge g_D(f)+1\). If its completed gap is zero,
the remaining partner set has size \(m\ge3\), contains an original
directed cycle, and gives
\(t_D(f)\ge g_D(f)+m\). In MIN this implies
\(\eta_D(f)\ge m+1\ge4\). Equality at residual four has the stated
necessary shape \(m=3,k=s=0,g_D=1,t_D=4\), with nonempty original
\(P(f)\) and a saturated full packet. The positive-indegree assertion
belongs to the \emph{terminal} partner graph; only cycle existence is
lifted to the original graph. These partner vertices are not thereby
shown to be actual failure sources or to have positive residual.

At strong Good equality with negative completed gap, the corrected
partner-path theorem applies instead: residual two permits only original
nongood-second partners, while residual three permits at most one far
partner, reached by an original arc from a nongood-second partner. That
far partner may be exceptional in the packet. Neither the all-second
case nor strict Good surplus has been excluded. Thus the five remaining
residual-seven partitions
\[
 (3,3,1),\ (3,2,2),\ (3,2,1,1),\
 (2,2,2,1),\ (2,2,1,1,1)
\]
are forced into negative completed gap, but remain unresolved.

\paragraph{The separate dense-boundary parameters.}
In the chapters restricted to \(n=2\delta+3\), \(\sigma\) denotes
target-capacity surplus and \(K=\{\sigma>0\}\) its core. It is not
\(\eta\), and a heavy target \(\sigma\ge2\) need not be an actual
bad-direction failure source. The general core conclusions are
\(|K|\ge6\) and \(z\ge h+8\), where \(z\) counts whole
minimum-degree vertices and \(h\) counts gap-two vertices.
The stronger \(z\ge h+10\) requires \(\delta\ge10\); the
\(h+11\) conclusion has an additional degree--residual condition,
for example \(\delta>40R+12\). The predecessor inequalities and
odd-deletion bounds retain their own hypotheses. In particular, a
retained minimum-degree condition is not asserted to hold for every
chosen odd set.

\paragraph{Corrections and evidence boundaries.}
The saved proofs do not use the rejected assertions that the external
second boundary of a target packet always lies in its incoming set,
that an exceptional far head cannot belong to the remainder, or that
the released-direction completion family contains the canonical family
as a set. The proved comparison is of maximum ordinary forward score.
At a fixed tournament the normalized Fisher losing density coincides
with the normalized reverse-expansion density; these do not provide
independent choices of weights. Three inward-priority families need not
be compatible. The enlarged convenient-tail closure preserves the
stated feed invariants but does not inherit original partner-cycle
lifting without another proof. Corrected or disproved shortcuts are
recorded only as scope controls, not as theorems.

The imported Tahara excess manuscript is attributed related work, not
an original result of this article. The archived local \(\delta=7\)
packages concern two specified branches and have solver-log validation;
their presence does not constitute a proof here of the whole
\(\delta=7\) theorem or a formal unsatisfiability certificate. Search
attainment is an upper bound on an extremal objective, and a timeout is
unresolved. Complete source-level coverage and the retained artifact
provenance are recorded in the accompanying inventory. The article and
supplement establish the stated partial results; they neither prove the
full conjecture nor exhibit a counterexample.
